% ============================================================================
%  TEMPLATE LaTeX — Rodapé Institucional UFMT
%  Disciplina : Algoritmos 1
%  Professora : Rafaela Souza Francisco
%  Curso      : Bacharelado em Ciência da Computação — UFMT, Cuiabá
% ============================================================================
\documentclass[12pt, a4paper]{article}

% --- Codificação e idioma ---------------------------------------------------
\usepackage[utf8]{inputenc}
\usepackage[T1]{fontenc}
% \usepackage[brazil]{babel}  % descomente se tiver texlive-lang-portuguese

% --- Fonte ------------------------------------------------------------------
% \usepackage{lmodern}        % descomente se disponível no seu sistema

% --- Layout de página -------------------------------------------------------
\usepackage[
  top=2.5cm,
  bottom=3.0cm,       % espaço extra para o rodapé
  left=2.5cm,
  right=2.5cm
]{geometry}

% --- Pacotes necessários ----------------------------------------------------
\usepackage{graphicx}          % inserir imagens
\usepackage[brazil]{babel}
\usepackage{fancyhdr}          % cabeçalho/rodapé customizado
\usepackage{xcolor}            % cores
\usepackage{hyperref}          % links clicáveis
\usepackage{tikz}              % desenhos (usado na linha decorativa)
\usepackage{amsmath, amssymb}  % símbolos matemáticos (útil para Algoritmos)
\usepackage{lastpage}          % referência à última página

\usetikzlibrary{shapes.geometric, arrows}
\tikzstyle{iniciofim} = [ellipse, draw, minimum width=3cm, minimum height=1cm, text centered]
\tikzstyle{processo} = [rectangle, draw, minimum width=5cm, minimum height=1cm, text centered]
\tikzstyle{seta} = [->, thick]
\tikzstyle{decisao} = [diamond, minimum width=3.5cm, minimum height=2cm, text centered, draw=black]

% --- Cores institucionais UFMT ---------------------------------------------
\definecolor{ufmtAzul}{RGB}{0, 51, 102}       % azul escuro institucional
\definecolor{ufmtVerde}{RGB}{0, 105, 62}       % verde institucional
\definecolor{ufmtCinza}{RGB}{100, 100, 100}    % cinza para texto secundário

% --- Dados do documento (EDITE AQUI) ----------------------------------------
\newcommand{\NomeAluno}{Thalles Eduardo Silva Mota}             % ← Seu nome
\newcommand{\Disciplina}{Algoritmos 1}
\newcommand{\Professora}{Profª. Rafaela Souza Francisco}
\newcommand{\Curso}{Bacharelado em Ciência da Computação}
\newcommand{\Instituicao}{UFMT — Universidade Federal de Mato Grosso}
\newcommand{\Semestre}{2026/1}
\renewcommand{\abstractname}{Resumo}
% ← Semestre atual
% Se tiver o logo em arquivo, defina o caminho abaixo:
% \newcommand{\LogoPath}{logo_ufmt.png}

% ============================================================================
%  CONFIGURAÇÃO DO RODAPÉ  (o coração do template)
% ============================================================================
\pagestyle{fancy}
\fancyhf{}                     % limpa cabeçalho e rodapé padrão

% --- Cabeçalho (opcional — descomente se quiser) ----------------------------
% \fancyhead[L]{\small\color{ufmtCinza}\Disciplina}
% \fancyhead[R]{\small\color{ufmtCinza}\NomeAluno}
% \renewcommand{\headrulewidth}{0.4pt}

% --- Rodapé -----------------------------------------------------------------
\renewcommand{\footrulewidth}{0pt}  % remove linha padrão (usaremos tikz)

\fancyfoot[C]{%
  \vspace{2pt}%
  % ── Linha decorativa com gradiente ──────────────────────────────────
  \begin{tikzpicture}[remember picture, overlay]
    \draw[
      line width=1.2pt,
      ufmtAzul
    ] ([yshift=8pt]current page.south west |- 0,0) ++(2.5cm,0)
      -- ++(\dimexpr\textwidth\relax, 0);
    % Pequeno detalhe em verde (identidade visual)
    \fill[ufmtVerde]
      ([yshift=6.5pt]current page.south west |- 0,0) ++(2.5cm,0)
      rectangle ++(4cm, 1.8pt);
  \end{tikzpicture}%
  %
  % ── Conteúdo do rodapé ──────────────────────────────────────────────
  \begin{minipage}[t]{0.12\textwidth}
    \vspace{-2pt}
    % =====================================================================
    % LOGO: Descomente a linha abaixo e comente o \fbox quando tiver o logo
    % \includegraphics[height=1.1cm]{\LogoPath}
    % =====================================================================
    % Placeholder visual enquanto não tiver o .png:
    \includegraphics[height=1.1cm]{images/UFMT - logo texto.png}%
  \end{minipage}%
  \hfill
  \begin{minipage}[t]{0.60\textwidth}
    \centering
    \vspace{-2pt}
    {\footnotesize\color{ufmtAzul}\textbf{\Instituicao}}\\[1pt]
    {\scriptsize\color{ufmtCinza}%
      \Curso\ \textbullet\ \Disciplina\ — \Semestre}\\[1pt]
    {\scriptsize\color{ufmtCinza}\Professora}
  \end{minipage}%
  \hfill
  \begin{minipage}[t]{0.15\textwidth}
    \raggedleft
    \vspace{-2pt}
    {\footnotesize\color{ufmtAzul}\textbf{\thepage}\,/\,\pageref{LastPage}}
  \end{minipage}%
}

% Estilo da primeira página (pode ser diferente)
\fancypagestyle{plain}{%
  \fancyhf{}%
  \renewcommand{\headrulewidth}{0pt}%
  \fancyfoot[C]{%
    \vspace{2pt}%
    \begin{tikzpicture}[remember picture, overlay]
      \draw[line width=1.2pt, ufmtAzul]
        ([yshift=8pt]current page.south west |- 0,0) ++(2.5cm,0)
        -- ++(\dimexpr\textwidth\relax, 0);
      \fill[ufmtVerde]
        ([yshift=6.5pt]current page.south west |- 0,0) ++(2.5cm,0)
        rectangle ++(4cm, 1.8pt);
    \end{tikzpicture}%
    \begin{minipage}[t]{0.12\textwidth}
      \vspace{-2pt}
      \fbox{\footnotesize\textbf{\color{ufmtAzul}UFMT}}%
    \end{minipage}%
    \hfill
    \begin{minipage}[t]{0.60\textwidth}
      \centering\vspace{-2pt}
      {\footnotesize\color{ufmtAzul}\textbf{\Instituicao}}\\[1pt]
      {\scriptsize\color{ufmtCinza}%
        \Curso\ \textbullet\ \Disciplina\ — \Semestre}\\[1pt]
      {\scriptsize\color{ufmtCinza}\Professora}
    \end{minipage}%
    \hfill
    \begin{minipage}[t]{0.15\textwidth}
      \raggedleft\vspace{-2pt}
      {\footnotesize\color{ufmtAzul}\textbf{\thepage}\,/\,\pageref{LastPage}}
    \end{minipage}%
  }%
}

% ============================================================================
%  DOCUMENTO DE EXEMPLO
% ============================================================================
\begin{document}

\begin{center}
  {\Large\bfseries\color{ufmtAzul} Tarefa: Algoritmos e Fluxogramas}\\[6pt]
  {\large \NomeAluno}\\[3pt]
  {\normalsize\color{ufmtCinza} \today}
\end{center}

\vspace{0.5cm}

\begin{abstract}
    \textbf{    Esta tarefa tem como objetivo apresentar e aplicar conceitos 
    fundamentais de algoritmos e lógica de programação por meio de quatro 
    problemas simples. Ao longo da tarefa, Serão exploradas estruturas condicionais com base na
análise de fluxogramas, permitindo a compreensão do processo de tomada
de decisão em algoritmos.}

\end{abstract}

\noindent\textbf{Palavras-chave: Algoritmos; Fluxogramas; Soluções Computacionais.} .

\section*{Introdução}
       Neste trabalho, serão abordados quatro exercícios que exploram conceitos 
    fundamentais para a criação de algoritmos e fluxogramas, permitindo a 
    compreensão prática de como construir soluções lógicas e eficientes para os
    mais diversos problemas.

\section*{Exercício 1 - Calcular o troco de uma compra}

\begin{center}
\begin{tikzpicture}[node distance=1.333cm]
 
% Nós
\node (inicio) [iniciofim] {Início};
\node (lerpreco) [processo, below of=inicio] {Ler preço do produto};
\node (lervalor) [processo, below of=lerpreco] {Ler valor entregue pelo cliente};
\node (suficiente) [decisao, below of=lervalor, yshift=-1.333cm] {Valor suficiente?};
\node (troco) [processo, below right of=suficiente, xshift=1.5cm, yshift=-1.5cm] {Troco $=$ valor $-$ preço};
\node (insuficiente) [processo, below left of=suficiente, xshift=-1.5cm, yshift=-1.5cm] {Valor insuficiente};
\node (exibir) [processo, below of=troco] {Exibir o troco};
\node (fim) [iniciofim, below of=suficiente, yshift=-3.8cm] {Fim};
 
% Setas
\draw [seta] (inicio) -- (lerpreco);
\draw [seta] (lerpreco) -- (lervalor);
\draw [seta] (lervalor) -- (suficiente);
\draw [seta] (suficiente.east) -- ++(1,0) node[above, midway] {Sim} -| (troco);
\draw [seta] (suficiente.west) -- ++(-1,0) node[above, midway] {Não} -| (insuficiente);
\draw [seta] (troco) -- (exibir);
\draw [seta] (exibir) |- (fim);
\draw [seta] (insuficiente) |- (fim);
 
\end{tikzpicture}
\end{center}

O primeiro exercício consiste em calcular o troco de uma compra, onde o 
    usuário deve informar o valor total da compra e o valor pago. 
    O algoritmo deve verificar se o valor pago é suficiente para cobrir a compra e, 
    caso seja, \textbf{calcular e exibir o troco}. Caso contrário, deve informar que o 
    valor é insuficiente(exemplificado na seção "Comandos de Entrada e Saída" do capítulo 2 de:\cite{forbellone2005}).

O algoritmo inicia com o passo \textbf{Início}. Em seguida, é realizada a 
leitura do \textbf{preço do produto}. Logo após, o programa solicita e lê 
o \textbf{valor entregue pelo cliente}.

    Com esses dados, o algoritmo realiza um cálculo simples: determina-se o 
\textbf{troco} por meio da subtração entre o valor entregue e o preço do 
produto, ou seja, \textbf{troco $=$ valor $-$ preço}. Após o cálculo, o resultado
é apresentado ao usuário por meio da exibição do \textbf{valor do troco}. 
Por fim, o algoritmo é encerrado no passo \textbf{Fim}.

    

\section*{Exercício 2 - Verificar se um número é par ou ímpar }

\begin{center}
\begin{tikzpicture}[node distance=1.5cm]
% Nós
\node (inicio) [iniciofim] {Início};
\node (entrada) [processo, below of=inicio] {Insira um número};
\node (decisao) [decisao, below of=entrada, yshift=-1cm] {NUM \% 2 = 0};
\node (impar) [processo, below left of=decisao, xshift=-3cm] {ÍMPAR};
\node (par) [processo, below right of=decisao, xshift=3cm] {PAR};
\node (fim) [iniciofim, below of=decisao, yshift=-1.5cm] {Fim};
% Setas
\draw [seta] (inicio) -- (entrada);
\draw [seta] (entrada) -- (decisao);
\draw [seta] (decisao) -| node[above, pos=0.25] {NÃO} (impar);
\draw [seta] (decisao) -| node[above, pos=0.25] {Sim} (par);
\draw [seta] (impar) |- (fim);
\draw [seta] (par) |- (fim);
\end{tikzpicture}
\end{center}

     O segundo exercício tem como objetivo verificar se um número é \textbf{par ou ímpar}.
     O usuário deve inserir um número inteiro, e o algoritmo deve determinar se 
     ele é \textbf{par (divisível por 2)} ou \textbf{ímpar (não divisível por 2)}, exibindo a 
     resposta correspondente a cada um dos casos.
     
    O algoritmo inicia com o passo \textbf{Início}. Em seguida, é realizada a leitura de um número inteiro informado pelo usuário. Após isso, o programa faz uma verificação: calcula-se o resto da \textbf{divisão desse número por 2}(exemplificado na seção 1.1 do capítulo 1 de:\cite{guidorizzi2001}).
    Se o resto da divisão for igual a zero, conclui-se que o número é \textbf{Par}, e essa informação é
    exibida ao usuário. Caso contrário, se o resto for diferente de zero, o número é considerado \textbf{Ímpar}, e essa mensagem é apresentada. Por fim, o algoritmo é encerrado no passo \textbf{Fim}.
\pagebreak

\section*{Exercício 3 - Encontrar o maior entre três números}

\begin{center}
\begin{tikzpicture}[node distance=1.333cm]

% Nós
\node (inicio) [iniciofim] {Início};
\node (ler) [processo, below of=inicio] {Ler A, B, C};
\node (ab) [decisao, below of=ler, yshift=-1cm] {A $>$ B?};
\node (ac) [decisao, below of=ab, xshift=-3cm, yshift=-1.5cm] {A $>$ C?};
\node (bc) [decisao, below of=ab, xshift=3cm, yshift=-1.5cm] {B $>$ C?};
\node (maiorA) [processo, below of=ac, xshift=-2.5cm, yshift=-1.5cm] {Maior $=$ A};
\node (maiorC1) [processo, below of=ac, xshift=2.5cm, yshift=-1.5cm] {Maior $=$ C};
\node (maiorB) [processo, below of=bc, xshift=-2.5cm, yshift=-3cm] {Maior $=$ B};
\node (maiorC2) [processo, below of=bc, xshift=2.5cm, yshift=-3cm] {Maior $=$ C};
\node (exibir) [processo, below of=ab, yshift=-8cm] {Exibir o maior};
\node (fim) [iniciofim, below of=exibir] {Fim};
 
% Setas principais
\draw [seta] (inicio) -- (ler);
\draw [seta] (ler) -- (ab);
 
% A > B?
\draw [seta] (ab.west) -- node[above] {Sim} (ac.north |- ab.west) -- (ac);
\draw [seta] (ab.east) -- node[above] {Não} (bc.north |- ab.east) -- (bc);
 
% A > C?
\draw [seta] (ac.west) -- node[above] {Sim} (maiorA.north |- ac.west) -- (maiorA);
\draw [seta] (ac.east) -- node[above] {Não} (maiorC1.north |- ac.east) -- (maiorC1);
 
% B > C?
\draw [seta] (bc.west) -- node[above] {Sim} (maiorB.north |- bc.west) -- (maiorB);
\draw [seta] (bc.east) -- node[above] {Não} (maiorC2.north |- bc.east) -- (maiorC2);
 
% Convergência para Exibir
\draw [seta] (maiorA) |- (exibir);
\draw [seta] (maiorC1) |- (exibir.north);
\draw [seta] (maiorB) |- (exibir.north);
\draw [seta] (maiorC2) |- (exibir);
 
\draw [seta] (exibir) -- (fim);
 
\end{tikzpicture}
\end{center}

O terceiro exercício tem como objetivo encontrar o \textbf{maior entre três números inteiros}. O usuário deve inserir \textbf{três números inteiros ($A, B, C$)}, e o algoritmo deve compará-los entre si para determinar qual deles possui o maior valor, exibindo a resposta correspondente(exemplificado na seção "Seleção Encadeada" do capítulo 3 de:
\cite{forbellone2005}). 

O algoritmo inicia com o passo \textbf{Início}. Em seguida, são realizadas as leituras de três números inteiros, informados pelo usuário, armazenados nas \textbf{variáveis ($A, B, C$)}.
Após a leitura dos valores, o programa inicia uma série de verificações utilizando decisões encadeadas. 

Primeiramente, verifica-se se o $A$ é maior que $B$ e também maior que $C$. 
Caso essa condição seja verdadeira, conclui-se que o $A$ é o maior, e essa informação é exibida ao usuário. Caso contrário, o algoritmo verifica se $B$ é maior que $A$ e também maior que $C$. Se essa condição for satisfeita, $B$ é considerado o maior, e o resultado é apresentado. Se nenhuma das condições anteriores for verdadeira, conclui-se que $C$ possui o maior valor entre os três, sendo então exibido ao usuário.
Por fim, o algoritmo é encerrado no passo \textbf{Fim}.

\section*{Exercício 4 - Contar de 1 até N}

\begin{center}
\begin{tikzpicture}[node distance=1.333cm]

% Nós
\node (inicio) [iniciofim] {Início};
\node (ler) [processo, below of=inicio] {Ler N};
\node (init) [processo, below of=ler] {i $\leftarrow$ 1};
\node (cond) [decisao, below of=init, yshift=-1cm] {i $\leq$ N?};
\node (exibir) [processo, below of=cond, xshift=3cm, yshift=-1.5cm] {Exibir i};
\node (incremento) [processo, below of=exibir] {i $\leftarrow$ i $+$ 1};
\node (fim) [iniciofim, below of=cond, xshift=-3cm, yshift=-1.5cm] {Fim};

% Setas principais
\draw [seta] (inicio) -- (ler);
\draw [seta] (ler) -- (init);
\draw [seta] (init) -- (cond);

% Decisão
\draw [seta] (cond.east) -- node[above] {Sim} (exibir.north |- cond.east) -- (exibir);
\draw [seta] (cond.west) -- node[above] {Não} (fim.north |- cond.west) -- (fim);

% Exibir → Incrementar
\draw [seta] (exibir) -- (incremento);

% Loop: Incrementar volta para a condição
\draw [seta] (incremento.south) -- ++(0,-0.8) -| (cond.south);

\end{tikzpicture}
\end{center}

    O quarto exercício tem como objetivo contar todos os \textbf{números inteiros de 1 até $N$}.
O usuário deve inserir um \textbf{número inteiro positivo $N$}, e o algoritmo deve apresentar
todos os valores de 1 até $N$, um por vez, utilizando uma estrutura de repetição(exemplificado na seção "Estruturas de Repetição" do capítulo 3 de:\cite{forbellone2005}).

O algoritmo inicia com o passo \textbf{Início}. Em seguida, é realizada a leitura de um número
inteiro positivo $N$, informado pelo usuário. Após isso, o programa inicializa uma \textbf{variável
chamada $i$} com o valor 1.

Em seguida, é feita uma verificação: enquanto o \textbf{valor de $i$ for menor ou igual a $N$}, o
algoritmo executa um conjunto de ações. Dentro dessa estrutura de repetição, o valor atual
de $i$ é exibido ao usuário. Logo após, a variável $i$ é incrementada em uma unidade.
Esse processo se repete continuamente até que \textbf{$i$ ultrapasse o valor de $N$}. Quando a
condição deixa de ser satisfeita, o algoritmo encerra sua execução no passo \textbf{Fim}.
\pagebreak

\section*{Conclusão}
       Portanto, perante os exercícios propostos, foi possível aplicar os conceitos de 
    algoritmos e fluxogramas para resolver problemas práticos, fortalecendo o 
    raciocínio lógico e a capacidade de organizar ideias na construção de 
    soluções computacionais. Além disso, os exercícios Contribuíram diretamente para o fortalecimento 
    do pensamento computacional, especialmente na tomada de decisões lógicas 
    e na execução de processos de forma sequencial e eficiente.
       Dessa forma, os exercícios cumpriram um papel importante na 
    consolidação dos conhecimentos iniciais em algoritmos, servindo como 
    base para o aprendizado de conceitos mais avançados na lógica de programação.

\bibliographystyle{abntex2-alf}  % ou plain, ieeetr, etc.
\bibliography{references}

\end{document}
